% Source: https://artofproblemsolving.com/community/c3h3180245_2019_abmc_online_november_contest__actonboxborough_math_competition
% Stable identity: aops-topic-3180245-q09
Each two-digit block with decimal digits $t,u$ has hexadecimal value $16t+u$. Since $16^2=256\equiv1\pmod{51}$, the whole number is congruent to the sum of its remaining block values. The prefix through 89 has sum $\sum_{t=1}^{8}\sum_{u=0}^{9}(16t+u)=160\cdot36+8\cdot45=6120=51\cdot120$, so removing the last ten blocks works. Prefixes ending 90 through 99 add the first $r=1,\ldots,10$ terms of 144,145,...,153. Their residues modulo 51 are 42, 34, 27, 21, 16, 12, 9, 7, 6, 6, none zero. Therefore no smaller deletion works and $\boxed{n=10}$.
